\section*{Bab \ref{ch:b2:catalytic-cycles} --- Katalisis Organologam: Tahap Elementer dan Siklus}

\begin{solution}{exo:b2:catalytic-cycles:1}
Sebelum: Ir(I), $d^8$, $8 + 4 \times 2 = 16$ elektron. Sesudah: Ir(III), $d^6$, $6 + 6 \times 2 = 18$
elektron, oktahedral.
\end{solution}

\begin{solution}{exo:b2:catalytic-cycles:2}
Yang pertama adalah \omterm{def:b2:catalytic-cycles:oxidative-addition}{eliminasi reduktif} (Pt(IV) menjadi Pt(II), etana
terbentuk dari dua metil). Yang kedua adalah \omterm{def:b2:catalytic-cycles:ligand-exchange}{pertukaran ligan} (satu fosfina
digantikan oleh etena) yang diikuti oleh \omterm{def:b2:catalytic-cycles:insertion}{insersi migrasi} etena ke dalam
ikatan Pd--Ph.
\end{solution}

\begin{solution}{exo:b2:catalytic-cycles:3}
TON $= 15.0/0.020 = 750$; TOF $= 750/3.0 = \qty{250}{h^{-1}}$.
\end{solution}

\begin{solution}{exo:b2:catalytic-cycles:4}
\begin{align*}
  &\mathrm{PdL_2 + ArBr \to ArPdBrL_2}\\
  &\mathrm{ArPdBrL_2 + Ar'B(OH)_3^- \to ArPdAr'L_2 + Br^- + B(OH)_3}\\
  &\mathrm{ArPdAr'L_2 \to Ar{-}Ar' + PdL_2}
\end{align*}
Jumlah: $\mathrm{ArBr + Ar'B(OH)_3^- \to Ar{-}Ar' + Br^- + B(OH)_3}$; setiap
spesies paladium saling meniadakan.
\end{solution}

\begin{solution}{exo:b2:catalytic-cycles:5}
Etil dan isopropil (keduanya mempunyai hidrogen $\beta$). Metil tidak mempunyai karbon $\beta$;
karbon $\beta$ neopentil tidak membawa hidrogen; karbon $\beta$ benzil adalah karbon cincin
\omterm{def:b2:huckel:aromatic}{aromatik} yang tidak membawa hidrogen.
\end{solution}

\begin{solution}{exo:b2:catalytic-cycles:6}
Metil (\textit{E})-3-fenilpropenoat, \ce{Ph-CH=CH-COOCH3}, dengan fenil yang
teradisi pada karbon ujung, produk trans.
\end{solution}

\begin{solution}{exo:b2:catalytic-cycles:7}
Butanal \ce{CH3CH2CH2CHO} (linear: logam beradisi pada karbon ujung, hidrida
pada karbon dalam) dan 2-metilpropanal \ce{(CH3)2CHCHO} (bercabang: logam
pada karbon dalam).
\end{solution}

\begin{solution}{exo:b2:catalytic-cycles:8}
\ce{[RhCl(PPh3)3]} mempunyai 16 elektron: adisi \ce{H2} (+2) akan
menghasilkan 18 dengan enam posisi koordinasi, terlalu sesak dengan tiga
fosfina yang meruah. \ce{[RhCl(PPh3)2]} mempunyai 14 dan sebuah situs
terbuka: \omterm{def:b2:catalytic-cycles:oxidative-addition}{adisi oksidatif} dengan mudah menghasilkan kompleks 16 elektron.
\end{solution}

\begin{solution}{exo:b2:catalytic-cycles:9}
Asam boronat adalah nukleofil yang buruk; basa menambahkan hidroksida (atau
alkoksida) pada boron dan menghasilkan borat \ce{ArB(OH)3-}, yang gugus
arilnya lebih kaya elektron dan berpindah ke paladium.
\end{solution}

\begin{solution}{exo:b2:catalytic-cycles:10}
Kompleks 20 elektron tidak masuk akal. Kompleks itu harus melepaskan sebuah
ligan lebih dahulu (14 e), lalu mengadisi \ce{H2} (16 e), mengikat alkena
(18 e), menyisipkannya ke dalam ikatan M--H (16 e), dan mengeliminasi alkana
(14 e), sebelum mengambil kembali ligannya atau memulai lagi.
\end{solution}

\begin{solution}{exo:b2:catalytic-cycles:11}
Laju awal $n_{\max}/\tau = \qty{0.25}{mmol/min}$; TOF awal $= 0.25/0.010 =
\qty{25}{min^{-1}} = \qty{1.5e3}{h^{-1}}$. TON akhir $= 10.0/0.010 = 1.0 \times 10^3$.
\end{solution}

\begin{solution}{exo:b2:catalytic-cycles:12}
4-Bromoanisol dengan asam fenilboronat (atau bromobenzena dengan asam
4-metoksifenilboronat), katalis fosfina paladium(0), dan sebuah basa.
Siklusnya: \omterm{def:b2:catalytic-cycles:oxidative-addition}{adisi oksidatif} aril bromida, \omterm{def:b2:catalytic-cycles:transmetalation}{transmetalasi} dari borat,
\omterm{def:b2:catalytic-cycles:oxidative-addition}{eliminasi reduktif} 4-metoksibifenil.
\end{solution}

\begin{solution}{pb:b2:catalytic-cycles:1}
\textbf{1.} Pd(II), $d^8$.
\textbf{2.} Pd(0), $d^{10}$, $10 + 2 \times 2 = 14$ elektron.
\textbf{3.} Dengan 14 elektron, kompleks itu dapat menerima dua pasang lagi:
kompleks itu mempunyai situs-situs kosong.
\textbf{4.} Kompleks fosfina paladium(0) dioksidasi oleh udara, dan fosfina
dioksidasi menjadi fosfina oksida.
\textbf{5.} Trifenilfosfina mereduksi paladium(II) menjadi paladium(0) (dan
dirinya sendiri teroksidasi) serta mengikat paladium(0), sehingga
mempertahankannya dalam larutan.
\textbf{6.} \omterm{def:b2:catalytic-cycles:cycle}{Prakatalis}.
\textbf{7.} \ce{Pd(PPh3)2 + CH3C6H4Br -> [Pd(C6H4CH3)Br(PPh3)2]}: Pd(II), $d^8$, $8 + 4 \times 2 =
16$ elektron.
\textbf{8.} Karbonat mengubahnya menjadi borat \ce{PhB(OH)3-}, yang
memindahkan gugus fenilnya.
\textbf{9.} \ce{[Pd(Ar)Br(PPh3)2] + PhB(OH)3- -> [Pd(Ar)(Ph)(PPh3)2] + Br- + B(OH)3}.
\textbf{10.} \omterm{def:b2:catalytic-cycles:oxidative-addition}{Eliminasi reduktif} menyambungkan dua ligan yang cis; aril-aril
yang trans harus berisomerisasi lebih dahulu.
\textbf{11.} \ce{[Pd(Ar)(Ph)(PPh3)2] -> Ar-Ph + Pd(PPh3)2}: paladium(0), 14
elektron, siap untuk putaran berikutnya.
\textbf{12.} \ce{CH3C6H4Br + PhB(OH)3- -> CH3C6H4-Ph + Br- + B(OH)3}.
\textbf{13.} 4-Bromotoluena (\qty{10.0}{mmol}, dibandingkan dengan \qty{12.0}{mmol} asam boronat).
\textbf{14.} $1.51/168.24 = \qty{8.98e-3}{mol}$, yaitu \qty{8.98}{mmol}.
\textbf{15.} $8.98/10.0 = 89.8~\%$.
\textbf{16.} $0.050/10.0 = 0.50$~mol~\%.
\textbf{17.} $8.98/0.050 = 180$.
\textbf{18.} $180/4.0 = \qty{45}{h^{-1}}$.
\textbf{19.} Dua gugus fenil berpindah ke paladium yang sama (homokopling
asam boronat, yang didorong oleh sekelumit oksigen), lalu dieliminasi
bersama-sama.
\textbf{20.} Gugus-gugus aril tidak mempunyai hidrogen $\beta$ yang dapat mencapai logam.
\textbf{21.} \omterm{def:b2:catalytic-cycles:oxidative-addition}{Adisi oksidatif}: ikatan \ce{C-Cl} lebih kuat daripada \ce{C-Br}.
\textbf{22.} Paladium berakhir sebagai partikel paladium(0) atau garam dalam
residu; logam itu mahal dan beracun, dan diambil kembali.
\textbf{23.} Sebagian asam boronat hilang dalam reaksi-reaksi samping
(homokopling, lepasnya boron dari cincin); kelebihan itu menjaga bromida
tetap menjadi pembatas.
\textbf{24.} Paladium berputar $\boldsymbol{\approx \num{1.8e2}}$ kali.
\end{solution}
